\section{Event selection and normalization}
\label{sec:selection-normalization}

The prompt search starts from selected $e^+e^-$ invariant-mass histograms. The
selection is designed to retain prompt radiative trident events and possible prompt
$\Ap\to e^+e^-$ decays while suppressing accidental track combinations and converted
wide-angle bremsstrahlung (WAB) backgrounds. Across the HPS running periods, the common
features are oppositely charged tracks, ECal cluster matching, timing consistency,
vertex-quality requirements, and momentum selections appropriate to the beam energy.
The GPR layer does not apply additional event-level cuts; it receives the selected
mass histogram, mass-resolution model, and normalization inputs for each run period.

\begin{table*}[t]
\caption{\label{tab:selection-normalization}Compact event-selection and normalization
summary. The full cut table and radiative-fraction provenance remain in the analysis
note; this table retains the information needed for the PRD-scale narrative.}
\begin{ruledtabular}
\begin{tabular}{llll}
Item & 2015 & 2016 & 2021 validation subset \\
Trigger & Pair & Pair1 & Single-cluster positron \\
Topology & Prompt $e^+e^-$ & Prompt $e^+e^-$ & Target-constrained $e^+e^-$ \\
Momentum & Prompt $p_{\mathrm{sum}}$ & Prompt $p_{\mathrm{sum}}$ & $p_{\mathrm{sum}}<2.8$ GeV \\
Resolution & Moller/scaled $\Ap$ & Moller/full-energy/smeared $\Ap$ & target-constrained model \\
Normalization & Radiative fraction & Radiative fraction & radiative-fraction inputs \\
Status & unblinded & 10\% subset & 1\% subset \\
\end{tabular}
\end{ruledtabular}
\end{table*}


The signal model is a Gaussian line shape whose width is fixed by the run-dependent
mass resolution. The 2015 and 2016 parameterizations follow the published and internal
engineering-run resolution studies based on Moller calibration, momentum-smearing
checks, and simulated $\Ap$ samples \cite{HPS2015DarkPhoton,HPS2016PromptLong,
HPS2015InternalNote,HPS2016BumpInternal}. The 2021 workflow uses the target-constrained
upgraded-detector resolution model, broadened by the target-constrained data/Monte
Carlo smearing scale documented in the analysis note. Figure~\ref{fig:mass-resolution}
shows the smooth curves propagated through the GPR signal templates.

\begin{figure*}[t]
\centering
\ifigure{0.88\textwidth}{mass_resolution_summary.png}
\caption{Mass-resolution parameterizations used by the current GPR workflow. The 2016
curve includes the high-mass continuation used to avoid an unphysical improvement of
the resolution outside the published prompt-search window. The 2021 curve includes the
target-constrained data/Monte Carlo broadening scale used by the upgraded-detector
analysis.}
\label{fig:mass-resolution}
\end{figure*}

The conversion from a fitted signal yield to $\eps^2$ uses the radiative fraction,
defined as the fraction of the selected prompt background attributed to radiative
trident production,
\begin{equation}
f_{\mathrm{rad},d}(m)
=
\frac{dN_{\gamma^\ast,d}/dm}
{dN_{\mathrm{tri},d}/dm+dN_{\mathrm{wab},d}/dm},
\label{eq:frad}
\end{equation}
where $d$ labels the data sample, ``tri'' denotes the selected trident background, and
``wab'' denotes WAB events that reconstruct as $e^+e^-$ candidates. The local prompt
mass density $\rho_d(m)$ is estimated from the observed selected spectrum in a
resolution-scaled mass interval. The yield-per-unit-coupling factor is
\begin{equation}
K_d(m)\equiv A_d(\eps^2=1)
=
\frac{3\pi\,m\,f_{\mathrm{rad},d}(m)\,\rho_d(m)}
{2\alpha_{\mathrm{EM}}},
\label{eq:kd}
\end{equation}
so a common dark-photon coupling predicts a data-set-dependent total signal yield
$A_d=K_d(m)\eps^2$. The current production configuration keeps the prompt-density
normalization window at $1.64\,\sigma_m$ even when the extraction likelihood uses the
wider $2.25\,\sigma_m$ signal window. This separates the historical prompt-density
normalization from the wider signal-recovery window used by the likelihood.

Radiative-fraction uncertainties are handled as normalization variations. When a
radiative penalty $\delta_{f,d}$ is enabled, the effective conversion is weakened by
$K_d^{\mathrm{eff}}=(1-\delta_{f,d})K_d$, so the same signal-yield limit corresponds
to a weaker limit on $\eps^2$. This penalty changes the coupling interpretation; it is
not part of the GPR background fit itself.
